\chapter{Derivation of the Expanded Welch--Satterthwaite Test for MI}
\label{chap:expanded}

The Expanded Welch test uses the estimated MI difference, standard error,
and statistic from Chapter~\ref{chap:baselines}, with a Student \(t\)
reference distribution. Its degrees of freedom depend on how much the
pointwise-MI variance estimates fluctuate between samples. This chapter
derives that sampling variability and uses it to calculate the degrees of
freedom. The name Expanded Welch refers to expanding the pointwise-MI
variance as the table probabilities change.

The argument begins with the moments needed for the degrees of freedom.
We then differentiate the variance directly, keeping the effects of
changing probability weights and changing pointwise-MI values separate.
The resulting observation-level sensitivity gives an approximation to
sampling error, from which the required moments follow.

\section{What the degrees of freedom require}
\label{sec:component-df-choice}

The squared standard error adds two estimated variance contributions,
\(\Vhat(P)/n_P\) and \(\Vhat(Q)/n_Q\). Under the
Welch--Satterthwaite principle, a contribution receives fewer degrees of
freedom when it varies more between samples relative to its mean. Fewer
degrees of freedom give the Student \(t\) distribution heavier tails
and a larger rejection cutoff, as described in
Section~\ref{sec:welch-background}.

To obtain the component degrees of freedom for \(P\), Satterthwaite
matches the sampling mean and variance of \(\Vhat(P)\) to those of a
variable with a scaled chi-squared distribution
\citep{satterthwaite1946}. For positive degrees of freedom
\(\nu_{P,\mathrm{mom}}\) to be chosen, write the approximating variable as
\begin{equation}
  \E\{\Vhat(P)\}
  \frac{\chi^2_{\nu_{P,\mathrm{mom}}}}{\nu_{P,\mathrm{mom}}},
  \label{eq:scaled-chi-working}
\end{equation}
A variable with this chi-squared distribution has mean
\(\nu_{P,\mathrm{mom}}\) and variance
\(2\nu_{P,\mathrm{mom}}\). The scaling therefore gives
\begin{align*}
  \E\left[
    \E\{\Vhat(P)\}
    \frac{\chi^2_{\nu_{P,\mathrm{mom}}}}{\nu_{P,\mathrm{mom}}}
  \right]
  &=\E\{\Vhat(P)\},\\
  \Var\left[
    \E\{\Vhat(P)\}
    \frac{\chi^2_{\nu_{P,\mathrm{mom}}}}{\nu_{P,\mathrm{mom}}}
  \right]
  &=\frac{[\E\{\Vhat(P)\}]^2}{\nu_{P,\mathrm{mom}}^2}
    (2\nu_{P,\mathrm{mom}})
   =\frac{2[\E\{\Vhat(P)\}]^2}{\nu_{P,\mathrm{mom}}}.
\end{align*}
The mean already matches. For positive sampling mean and variance,
equating the last expression to \(\Var\{\Vhat(P)\}\) and solving
gives the component moment target:
\begin{equation}
  \boxed{
  \nu_{P,\mathrm{mom}}
  =\frac{2[\E\{\Vhat(P)\}]^2}
         {\Var\{\Vhat(P)\}}
  }.
  \label{eq:component-df-moments}
\end{equation}
Dividing \(\Vhat(P)\) by the fixed sample size \(n_P\) scales both
its squared mean and its variance by \(1/n_P^2\), so their ratio and
the component degrees of freedom are unchanged. The same calculation
applies to \(Q\).

We therefore need the sampling mean and variance of each pointwise-MI
variance estimate. We obtain them by calculating how one observation
affects the estimate, then averaging these contributions over an
independent sample.

\section{How pointwise-MI variance changes with the table}
\label{sec:variance-sensitivity}

For a positive population table, write the variance as a weighted sum of
squared deviations from MI:
\begin{equation}
  V(P)=\sum_{i,j}p_{ij}\{\pmi_P(i,j)-I(P)\}^2.
  \label{eq:v-centred-form}
\end{equation}
A change in the table affects both the probability weights \(p_{ij}\)
and the pointwise-MI values \(\pmi_P(i,j)\). Differentiating this
centred expression lets us calculate these effects separately and add them.

\subsection{A table change representing one observation}

An observation in cell \((x,y)\) can be represented by a table
\(\delta_{(x,y)}\) that places probability one in that cell. To
measure its effect, consider a small movement from the population table
towards this one-observation table
\citep{hampel1974,vanderVaart1998}:
\begin{equation}
  P_\varepsilon=(1-\varepsilon)P+\varepsilon\delta_{(x,y)}.
  \label{eq:contamination-path}
\end{equation}
At \(\varepsilon=0\), the table is \(P\). Increasing
\(\varepsilon\) reduces the original probabilities proportionally and
adds probability to \((x,y)\), keeping total probability one. The
empirical table is an average of one-observation tables, so this path
isolates one observation's contribution to its difference from \(P\).

Define the observation-level sensitivity, or influence function, by
\begin{equation}
  g_P(x,y)
  =\left.\frac{\mathrm d V(P_\varepsilon)}{\mathrm d\varepsilon}
   \right|_0.
  \label{eq:g-definition}
\end{equation}
This derivative measures the change in \(V\) per unit increase in the
mixture weight at the original table. We evaluate it using the changes in
the cell probabilities and their row and column margins.

The selected cell \((x,y)\) stays fixed, while \((i,j)\) indexes
any cell whose probability is being evaluated. Along the path,
\begin{equation}
  p_{ij}(\varepsilon)
  =(1-\varepsilon)p_{ij}+\varepsilon\mathbf 1\{i=x,j=y\}.
\end{equation}
The joint, row, and column probability derivatives at zero are
\begin{align}
  \left.\frac{\mathrm d p_{ij}(\varepsilon)}{\mathrm d\varepsilon}
  \right|_0
  &=\mathbf 1\{i=x,j=y\}-p_{ij},
  \label{eq:joint-perturbation-derivative}\\
  \left.\frac{\mathrm d p_{i+}(\varepsilon)}{\mathrm d\varepsilon}
  \right|_0
  &=\mathbf 1\{i=x\}-p_{i+},\\
  \left.\frac{\mathrm d p_{+j}(\varepsilon)}{\mathrm d\varepsilon}
  \right|_0
  &=\mathbf 1\{j=y\}-p_{+j}.
\end{align}

\Needspace{24\baselineskip}
\subsection{Differentiating the centred variance}

Along the path, Equation~\eqref{eq:v-centred-form} becomes
\begin{equation*}
  V(P_\varepsilon)
  =\sum_{i,j}p_{ij}(\varepsilon)
    \{\pmi_{P_\varepsilon}(i,j)-I(P_\varepsilon)\}^2.
\end{equation*}
The product rule differentiates the probability weights and the squared
deviations. Within each squared deviation, both pointwise MI and its
mean can change. Evaluating at zero gives
\begin{align}
  g_P(x,y)
  ={}&\sum_{i,j}\{\mathbf 1(i=x,j=y)-p_{ij}\}
      \{\pmi_P(i,j)-I(P)\}^2\notag\\
  &+2\sum_{i,j}p_{ij}\{\pmi_P(i,j)-I(P)\}
    \left[
      \left.\frac{\mathrm d\pmi_{P_\varepsilon}(i,j)}
                       {\mathrm d\varepsilon}\right|_0
      -\left.\frac{\mathrm d I(P_\varepsilon)}
                       {\mathrm d\varepsilon}\right|_0
    \right].
  \label{eq:centred-variance-product-rule}
\end{align}
The derivative of \(I(P_\varepsilon)\) is the same for every cell,
so its coefficient is the probability-weighted sum of centred
pointwise-MI values. That sum is zero:
\begin{equation*}
  \sum_{i,j}p_{ij}\{\pmi_P(i,j)-I(P)\}
  =\underbrace{\sum_{i,j}p_{ij}\pmi_P(i,j)}_{I(P)}
   -I(P)\underbrace{\sum_{i,j}p_{ij}}_1=0.
\end{equation*}
The mean-derivative term therefore contributes zero without requiring
its value to be calculated separately.

In the first sum of Equation~\eqref{eq:centred-variance-product-rule},
the indicator selects the squared deviation in cell \((x,y)\), and
the remaining probability-weighted sum is \(V(P)\). Thus
\begin{align*}
  &\sum_{i,j}\{\mathbf 1(i=x,j=y)-p_{ij}\}
      \{\pmi_P(i,j)-I(P)\}^2\\
  &\qquad=\{\pmi_P(x,y)-I(P)\}^2-V(P).
\end{align*}
Substituting these two simplifications leaves the two contributions
we need:
\begin{align}
  g_P(x,y)
  ={}&\{\pmi_P(x,y)-I(P)\}^2-V(P)\notag\\
  &+2\sum_{i,j}p_{ij}\{\pmi_P(i,j)-I(P)\}
    \left.\frac{\mathrm d\pmi_{P_\varepsilon}(i,j)}
                     {\mathrm d\varepsilon}\right|_0.
  \label{eq:g-direct-decomposition}
\end{align}
The first line is the effect of changing the probability weights with
the pointwise-MI values held fixed. The second line accounts for changes
in those values, which depend on the joint and marginal probabilities.

\subsection{The effect of changing pointwise MI}

Pointwise MI along the path is
\begin{equation*}
  \pmi_{P_\varepsilon}(i,j)
  =\log p_{ij}(\varepsilon)
   -\log p_{i+}(\varepsilon)
   -\log p_{+j}(\varepsilon).
\end{equation*}
Each logarithm has derivative equal to the probability derivative
divided by the probability. Substituting the derivatives above gives
\begin{align}
  \left.\frac{\mathrm d\pmi_{P_\varepsilon}(i,j)}
                   {\mathrm d\varepsilon}\right|_0
  &={\frac{\mathbf 1\{i=x,j=y\}-p_{ij}}{p_{ij}}}
    -\frac{\mathbf 1\{i=x\}-p_{i+}}{p_{i+}}
    -\frac{\mathbf 1\{j=y\}-p_{+j}}{p_{+j}}\notag\\
  &=\frac{\mathbf 1\{i=x,j=y\}}{p_{ij}}
    -\frac{\mathbf 1\{i=x\}}{p_{i+}}
    -\frac{\mathbf 1\{j=y\}}{p_{+j}}+1.
  \label{eq:pmi-derivative}
\end{align}
The constant is \(-1+1+1=1\): the joint logarithm contributes
\(-1\), and subtracting each marginal logarithm contributes \(+1\).

Insert this derivative into the sum in the second line of
Equation~\eqref{eq:g-direct-decomposition}. The indicators select the
chosen cell, its row, and its column:
\begin{align}
  &\sum_{i,j}p_{ij}\{\pmi_P(i,j)-I(P)\}
    \left.\frac{\mathrm d\pmi_{P_\varepsilon}(i,j)}
                     {\mathrm d\varepsilon}\right|_0\notag\\
  &\quad=\pmi_P(x,y)-I(P)
    -\frac{\sum_jp_{xj}\{\pmi_P(x,j)-I(P)\}}{p_{x+}}\notag\\
  &\qquad\quad
    -\frac{\sum_ip_{iy}\{\pmi_P(i,y)-I(P)\}}{p_{+y}}
    +\underbrace{\sum_{i,j}p_{ij}\{\pmi_P(i,j)-I(P)\}}_0.
  \label{eq:centred-score-sum}
\end{align}
For the row term, the indicator fixes \(i=x\), leaving the sum over
columns. Dividing by \(p_{x+}\) gives conditional probabilities
within that row, so
\begin{align*}
  \frac{\sum_jp_{xj}\{\pmi_P(x,j)-I(P)\}}{p_{x+}}
  &=\sum_j\frac{p_{xj}}{p_{x+}}\pmi_P(x,j)-I(P)\\
  &=\E_P\{\pmi_P(X,Y)\mid X=x\}-I(P).
\end{align*}
\begin{samepage}
The weights \(p_{xj}/p_{x+}\) sum to one, which gives the single
subtraction of \(I(P)\). The column term similarly equals
\(\E_P\{\pmi_P(X,Y)\mid Y=y\}-I(P)\), and the final sum is zero
by centring. Subtracting the row and column terms returns two copies of
\(I(P)\), while the selected-cell term subtracts one, leaving
\(+I(P)\). Substituting in Equation~\eqref{eq:g-direct-decomposition},
with its factor \(2\) on the pointwise-MI contribution, gives
\begin{equation}
\boxed{
\begin{aligned}
g_P(x,y)
={}&\{\pmi_P(x,y)-I(P)\}^2-V(P)\\
&+2\bigl[
\pmi_P(x,y)
-\E_P\{\pmi_P(X,Y)\mid X=x\}\\
&\hspace{18mm}
-\E_P\{\pmi_P(X,Y)\mid Y=y\}
+I(P)
\bigr].
\end{aligned}
}
\label{eq:g-final}
\end{equation}
The two contributions remain separate in the result. The first line
measures the effect of changing probability weights; the bracket measures
the effect of changing pointwise MI through the joint, row, and column
probabilities. Together they describe the sensitivity of \(V\) to an
observation in cell \((x,y)\).
\par
\end{samepage}

\section{From sensitivity to degrees of freedom}
\label{sec:component-df}

\subsection{Sampling error of the variance estimate}

The empirical table averages one-observation tables. If
\(Z_a^{(P)}\) is observation \(a\) from \(P\), then
\begin{equation}
  \widehat P-P
  =\frac{1}{n_P}\sum_{a=1}^{n_P}
    \{\delta_{Z_a^{(P)}}-P\}.
\end{equation}
Since \(\Vhat(P)=V(\widehat P)\), we obtain its first-order sampling
error by applying a Taylor expansion to this table change.

For a differentiable scalar function, the first-order formula at zero is
\begin{equation*}
  f(t)\approx f(0)+t f'(0).
\end{equation*}
Here set \(f(t)=V\{P+t(\widehat P-P)\}\): the expansion point is
the population table \(P\), the new table is \(\widehat P\), and
their difference is \(\widehat P-P\). The scalar \(t\) moves along
the line between them, with \(f(0)=V(P)\) and
\(f(1)=\Vhat(P)\). Evaluating at \(t=1\) gives
\begin{equation*}
  \Vhat(P)-V(P)
  \approx\left.\frac{\mathrm d}{\mathrm dt}
    V\{P+t(\widehat P-P)\}\right|_{t=0}.
\end{equation*}
The approximation requires the cell-probability errors
\(\widehat P-P\) to be small. The omitted terms are quadratic or
higher order in those errors, even though the endpoint of the path is
\(t=1\).

At fixed \(P\), the first-order Taylor term is linear in the table
change: each cell-probability change is multiplied by a coefficient
evaluated at \(P\). The derivative of the average table change is
therefore the average of the individual derivatives
\citep[Chapter~20]{vanderVaart1998}. Each direction
\(\delta_{Z_a^{(P)}}-P\) is the one-observation path used to define
\(g_P\), giving
\begin{equation}
  \Vhat(P)-V(P)
  \approx\frac{1}{n_P}\sum_{a=1}^{n_P}g_P(Z_a^{(P)}).
  \label{eq:vhat-linearization}
\end{equation}
Each observation thus contributes \(g_P(Z_a^{(P)})/n_P\) to the
first-order error. We can now calculate the mean and variance of this
average.

\subsection{Sampling mean and variance of the estimate}

The sensitivity has mean zero. For the first line of
Equation~\eqref{eq:g-final}, averaging the squared deviation gives
\(V(P)\), which cancels the subtracted \(V(P)\). For the bracket,
averaging each conditional mean recovers overall MI. For example,
\begin{align*}
  \E_P\left[\E_P\{\pmi_P(X,Y)\mid X\}\right]
  &=\sum_xp_{x+}\,
    \frac{\sum_yp_{xy}\pmi_P(x,y)}{p_{x+}}\\
  &=\sum_{x,y}p_{xy}\pmi_P(x,y)=I(P).
\end{align*}
The row probabilities cancel the conditional-mean denominators,
recovering the sum over all cells. The column calculation gives the same
result. Thus
\begin{equation}
  \E_P\{g_P(X,Y)\}
  =\{V(P)-V(P)\}
   +2\{I(P)-I(P)-I(P)+I(P)\}=0.
\end{equation}
Taking expectations in Equation~\eqref{eq:vhat-linearization} therefore
gives the first required moment:
\begin{equation}
  \E\{\Vhat(P)\}\approx V(P).
  \label{eq:vhat-mean-approximation}
\end{equation}
Bias in \(\Vhat(P)\) arises from terms beyond this linear approximation.

To calculate the sampling variance of \(\Vhat(P)\), define the
population variance of the sensitivities by
\begin{equation}
  \tau^2(P)=\Var_P\{g_P(X,Y)\}.
  \label{eq:tau-definition}
\end{equation}
Independent observations give independent sensitivity contributions, so
their covariance terms vanish. The sampling variance of their average is
the individual variance divided by \(n_P\):
\begin{align}
  \Var\{\Vhat(P)\}
  &\approx\frac{1}{n_P^2}
    \sum_{a=1}^{n_P}\Var_P\{g_P(Z_a^{(P)})\}\notag\\
  &=\frac{\tau^2(P)}{n_P}.
  \label{eq:vhat-sampling-variance}
\end{align}
These two approximations supply the sampling mean and variance needed
for moment matching.


\Needspace{12\baselineskip}
\subsection{Population and estimated component degrees of freedom}

The mean and variance approximations now supply both inputs to
Equation~\eqref{eq:component-df-moments}. To obtain positive, finite
component degrees of freedom, we require positive sensitivity variances:
\begin{assumption}[Positive variances of observation-level sensitivities]
For the first-order construction of the Expanded Welch test,
\(\tau^2(P)>0\) and \(\tau^2(Q)>0\).
\end{assumption}
Substituting the first-order mean and variance approximations into
Equation~\eqref{eq:component-df-moments} gives
\begin{equation*}
  \nu_{P,\mathrm{mom}}
  \approx\frac{2V(P)^2}{\tau^2(P)/n_P}
  =\frac{2n_PV(P)^2}{\tau^2(P)}.
\end{equation*}
We define this population first-order approximation by
\begin{equation}
  \boxed{
  \nu_{P,\mathrm{FO}}
  =
  \frac{2n_PV(P)^2}{\tau^2(P)}
  }.
  \label{eq:population-component-df}
\end{equation}
The ratio \(\tau^2(P)/\{n_PV(P)^2\}\) is the first-order sampling
variance of \(\Vhat(P)\) relative to \(V(P)^2\). The component degrees
of freedom are twice its reciprocal. Small relative fluctuations give
large degrees of freedom and a Student \(t\) distribution close to the
standard normal distribution; larger relative fluctuations give fewer
degrees of freedom and a larger rejection cutoff.
Positive \(V(P)\) does not guarantee positive \(\tau^2(P)\): table changes
can have no first-order effect on \(V(P)\). In that case, describing the
sampling variation of \(\Vhat(P)\) requires a higher-order approximation.

\medskip
\thesisminiheading{Table-based component degrees of freedom.}

The population inputs \(V(P)\) and \(\tau^2(P)\) are unknown, so the
test estimates them from the observed table. We replace every population
quantity in Equation~\eqref{eq:g-final} by its empirical counterpart and sum over
occupied cells, since pointwise MI is undefined at an empirical zero.
For each occupied cell,
\begin{equation}
\begin{aligned}
\widehat g_P(i,j)
={}&\{\widehat\pmi_P(i,j)-\Ihat(P)\}^2-\Vhat(P)\\
&+2\left[
\widehat\pmi_P(i,j)
-\frac{\sum_{j'}\widehat p_{ij'}\widehat\pmi_P(i,j')}
       {\widehat p_{i+}}
-\frac{\sum_{i'}\widehat p_{i'j}\widehat\pmi_P(i',j)}
       {\widehat p_{+j}}
+\Ihat(P)
\right].
\end{aligned}
\label{eq:g-hat}
\end{equation}
An empty cell has zero weight in these sums, so no logarithm or sensitivity
needs to be evaluated there. Although the empirical sensitivities have
probability-weighted mean zero under exact substitution, their variance
is centred explicitly:
\begin{equation}
  \widehat\tau^2(P)
  =\sum_{i,j}\widehat p_{ij}
  \left\{\widehat g_P(i,j)
  -\sum_{k,l}\widehat p_{kl}\widehat g_P(k,l)\right\}^2.
  \label{eq:tau-hat}
\end{equation}
The weighted variance \(\widehat\tau^2(P)\) estimates the population
variance of the observation-level sensitivities. By
Equation~\eqref{eq:vhat-sampling-variance}, \(\widehat\tau^2(P)/n_P\)
then estimates the sampling variance of \(\Vhat(P)\). Replacing \(V(P)\) and
\(\tau^2(P)\) by their estimates in
Equation~\eqref{eq:population-component-df} gives the estimated component
degrees of freedom:
\begin{equation}
  \boxed{
  \widehat\nu_P
  =\frac{2n_P\Vhat(P)^2}{\widehat\tau^2(P)}
  }.
  \label{eq:nu-p-hat}
\end{equation}
Repeating the same calculation for the second table gives
\(\widehat\nu_Q\).

\medskip
Table~\ref{tab:component-df-roles} distinguishes the sampling-moment target,
its population first-order approximation, and the sample-based estimate.
The target and population approximation are fixed for a population and
sample size, whereas \(\widehat\nu_P\) changes with each sampled table.

\begin{table}[htbp]
\centering
\caption{Component degrees of freedom and their inputs.}
\label{tab:component-df-roles}
\thesistableformat
\begin{tabularx}{\textwidth}{@{}l>{\raggedright\arraybackslash}X>{\raggedright\arraybackslash}X@{}}
\toprule
\textbf{Quantity} & \textbf{Role} & \textbf{Inputs}\\
\midrule
\(\nu_{P,\mathrm{mom}}\) & Sampling-moment target &
\(\E\{\Vhat(P)\}\) and \(\Var\{\Vhat(P)\}\) at fixed \(P,n_P\)\\
\addlinespace
\(\nu_{P,\mathrm{FO}}\) & Population first-order approximation to the target &
Population \(V(P)\), \(\tau^2(P)\), and \(n_P\)\\
\addlinespace
\(\widehat\nu_P\) & Component estimate used by Expanded Welch &
\(\Vhat(P)\), \(\widehat\tau^2(P)\), and \(n_P\) from one sample\\
\bottomrule
\end{tabularx}
\end{table}

\subsection{Combined degrees of freedom and p-value}

The two component degrees of freedom must now be combined to choose one
reference distribution for \(T\). For component inputs \(\nu_P\) and
\(\nu_Q\), the Welch--Satterthwaite equation is
\citep{satterthwaite1946,welch1947}:
\begin{equation}
  \nu
  =
  \frac{
  \left\{\Vhat(P)/n_P+\Vhat(Q)/n_Q\right\}^2
  }{
  \left\{\Vhat(P)/n_P\right\}^2/\nu_P
  +\left\{\Vhat(Q)/n_Q\right\}^2/\nu_Q
  }.
  \label{eq:ws-general}
\end{equation}
To derive this rule, we apply moment matching to the full estimated
variance,
\begin{equation*}
  \SEhat^2=\frac{\Vhat(P)}{n_P}+\frac{\Vhat(Q)}{n_Q}.
\end{equation*}
Linearity of expectation gives its sampling mean:
\begin{equation*}
  \E(\SEhat^2)
  =\frac{\E\{\Vhat(P)\}}{n_P}
   +\frac{\E\{\Vhat(Q)\}}{n_Q}.
\end{equation*}
The two pointwise-MI variance estimates are computed from independent
samples, so their covariance is zero and their contributions to
\(\Var(\SEhat^2)\) add. Rearranging
Equation~\eqref{eq:component-df-moments} expresses each component's sampling
variance in terms of its sampling mean and moment-matched degrees of
freedom. Applying this identity to both components gives
\begin{align*}
  \Var(\SEhat^2)
  &=\frac{\Var\{\Vhat(P)\}}{n_P^2}
    +\frac{\Var\{\Vhat(Q)\}}{n_Q^2}\\*
  &=2\left\{
    \frac{[\E\{\Vhat(P)\}/n_P]^2}{\nu_{P,\mathrm{mom}}}
    +\frac{[\E\{\Vhat(Q)\}/n_Q]^2}{\nu_{Q,\mathrm{mom}}}
    \right\}.
\end{align*}
The combined moment-matched degrees of freedom are twice the squared mean
of \(\SEhat^2\) divided by its variance. Substituting the expressions above
and cancelling the common factor \(2\) gives
\begin{equation*}
  \frac{2[\E(\SEhat^2)]^2}{\Var(\SEhat^2)}
  =\frac{
    \left\{\E\{\Vhat(P)\}/n_P+\E\{\Vhat(Q)\}/n_Q\right\}^2
    }{
    [\E\{\Vhat(P)\}/n_P]^2/\nu_{P,\mathrm{mom}}
    +[\E\{\Vhat(Q)\}/n_Q]^2/\nu_{Q,\mathrm{mom}}
    }.
\end{equation*}
This establishes the combining rule in Equation~\eqref{eq:ws-general} for
population moment inputs.

\begin{samepage}
For the observed tables, we replace the sampling
means \(\E\{\Vhat(P)\}\) and \(\E\{\Vhat(Q)\}\) by \(\Vhat(P)\) and
\(\Vhat(Q)\), and the component targets \(\nu_{P,\mathrm{mom}}\) and
\(\nu_{Q,\mathrm{mom}}\) by \(\widehat\nu_P\) and \(\widehat\nu_Q\).
These substitutions give the estimated combined degrees of freedom:
\begin{equation}
  \boxed{
  \widehat\nu_{\mathrm{expanded}}
  =
  \frac{
  \left\{\Vhat(P)/n_P+\Vhat(Q)/n_Q\right\}^2
  }{
  \left\{\Vhat(P)/n_P\right\}^2/\widehat\nu_P
  +
  \left\{\Vhat(Q)/n_Q\right\}^2/\widehat\nu_Q
  }
  }.
  \label{eq:expanded-df}
\end{equation}
\par
\end{samepage}
The two-sided p-value for the Expanded Welch test is
\begin{equation}
  p_{\mathrm{expanded}}
  =2\left\{1-F_t(|T|;\widehat\nu_{\mathrm{expanded}})\right\},
  \label{eq:expanded-p}
\end{equation}
where \(F_t(x;\nu)\) denotes the cumulative distribution function of a
Student \(t\) distribution with \(\nu\) degrees of freedom, evaluated at \(x\).

For a fixed finite \(T\), the two-sided p-value from the Student \(t\)
distribution is at least as large as the Wald p-value from the standard
normal distribution.
We call a test calculation successful when all required quantities satisfy
the numerical conditions in Appendix~\ref{sec:validity-rules}. Every
successful calculation for the Expanded Welch test also satisfies the Wald test's
calculation requirements. At the same p-value threshold, every rejection by the
Expanded Welch test is therefore also a rejection by the Wald test, including when unsuccessful calculations
count as non-rejections in the reported rates. At the stated significance
level, the Expanded Welch test cannot raise rejection where the Wald test is already
conservative.

\section{Worked example and implementation}

\subsection{Worked comparison}

Consider the two observed \(3\times3\) tables
\begin{equation}
  N^P=
  \begin{pmatrix}
  35&10&5\\
  10&20&5\\
  5&5&10
  \end{pmatrix},
  \qquad
  N^Q=
  \begin{pmatrix}
  30&5&5\\
  5&30&5\\
  5&5&10
  \end{pmatrix}.
  \label{eq:worked-tables}
\end{equation}
Their totals are \(n_P=105\) and \(n_Q=100\). In \(N^P\), cell
\((1,1)\) has count 35, and its row and column each total 50. Its empirical
pointwise MI is therefore
\begin{equation*}
 \widehat\pmi_P(1,1)
 =\log\left\{\frac{35/105}{(50/105)(50/105)}\right\}
 =\log(1.47)\simeq0.38526.
\end{equation*}
This cell contributes \((35/105)(0.38526)\) to the plug-in MI estimate.
Summing the contributions from all nine cells gives
\(\Ihat(P)=0.13693\); the same calculation for \(N^Q\) gives
\(\Ihat(Q)=0.25848\) nats. Since \((3-1)(3-1)=4\), the corrected difference is
\begin{equation}
  \Deltahat
  =\left(0.13693-\frac{4}{2(105)}\right)
  -\left(0.25848-\frac{4}{2(100)}\right)
  =-0.12059.
\end{equation}

The variance calculation weights each squared deviation from the estimated
MI by its cell probability. Cell \((1,1)\), for example, contributes
\((35/105)(0.38526-0.13693)^2\simeq0.02056\) to \(\Vhat(P)\).
Summing over all cells gives \(\Vhat(P)=0.27189\) and
\(\Vhat(Q)=0.43371\). Hence
\begin{equation}
  \SEhat
  =\left(\frac{0.27189}{105}+\frac{0.43371}{100}\right)^{1/2}
  =0.08323,
  \qquad
  T=-1.44898.
\end{equation}
The Wald test gives \(p=0.14734\).

For the Expanded Welch test, the same cell illustrates the sensitivity calculation.
The first row's pointwise-MI values are 0.38526, \(-0.51083\), and
\(-0.64436\), weighted by its counts 35, 10, and 5. Its conditional mean is
\begin{equation*}
 \frac{35(0.38526)+10(-0.51083)+5(-0.64436)}{50}
 \simeq0.10308.
\end{equation*}
Because this count table is symmetric, the first column has the same
conditional mean. Substituting these values into Equation~\eqref{eq:g-hat} gives
\begin{align*}
 \widehat g_P(1,1)
 &\simeq (0.38526-0.13693)^2-0.27189\\*
 &\quad+2\{0.38526-0.10308-0.10308+0.13693\}\\*
 &\simeq0.4218.
\end{align*}
The probability-weighted mean of the sensitivities is zero, so this cell
contributes \((35/105)(0.4218)^2\simeq0.0593\) to
\(\widehat\tau^2(P)\). Repeating the calculation for all cells in both
tables gives \(\widehat\tau^2(P)=0.93212\) and
\(\widehat\tau^2(Q)=0.82035\). Substitution into
Equation~\eqref{eq:nu-p-hat} gives component degrees of freedom
\begin{equation}
  \widehat\nu_P=16.65,
  \qquad
  \widehat\nu_Q=45.86.
\end{equation}
Combining them with Equation~\eqref{eq:expanded-df} gives
\(\widehat\nu_{\mathrm{expanded}}=59.03\) and \(p=0.15264\).
Neither method rejects at the 5\% level.

\pagebreak[3]
\subsection{Operational algorithm and complexity}
\label{sec:expanded-complexity}

Table~\ref{tab:expanded-algorithm} summarises one comparison.

\begin{table}[htbp]
\centering
\caption{Operational Expanded Welch--Satterthwaite algorithm.}
\label{tab:expanded-algorithm}
\thesistableformat
\begin{tabularx}{\textwidth}{@{}c>{\raggedright\arraybackslash}X@{}}
\toprule
\textbf{Step} & \textbf{Calculation} \\
\midrule
1 & Validate the two aligned count tables and their specified numbers of row and column categories; independence of samples is a design assumption. \\
\addlinespace[0.4em]
2 & Compute empirical probabilities, pointwise MI, plug-in MI, and the leading bias-corrected difference. \\
\addlinespace[0.4em]
3 & Compute \(\Vhat(P)\), \(\Vhat(Q)\), the shared standard error, and \(T\). \\
\addlinespace[0.4em]
4 & Compute the row and column pointwise-MI means needed in \(\widehat g_P\) and \(\widehat g_Q\). \\
\addlinespace[0.4em]
5 & Compute \(\widehat\tau^2(P)\), \(\widehat\tau^2(Q)\), and the two component degrees of freedom. \\
\addlinespace[0.4em]
6 & Combine the components using Equation~\eqref{eq:expanded-df}. \\
\addlinespace[0.4em]
7 & Return the p-value from the Student \(t\) distribution, effective degrees of freedom, and results of the calculation checks. \\
\bottomrule
\end{tabularx}
\end{table}

Each empirical probability, pointwise-MI value, row mean, column mean, and
variance contribution is calculated by a fixed number of table scans. The
time complexity is therefore \(O(rc)\), and the working memory is also
\(O(rc)\).

\section{Behaviour and limitations}
\label{sec:near-independence-mechanism}

\subsection{Near independence}

Consider a population table that approaches independence while its row
and column margins stay fixed. At independence, every pointwise-MI value
is zero. Since the linear MI error in
Equation~\eqref{eq:mi-error-from-cell-errors} weights each cell-probability
error by its population pointwise MI, this linear contribution vanishes.
Sampling errors can still create apparent association through products of
cell-probability errors, which are retained by a second-order expansion.
Near independence, the small linear weights allow these quadratic terms
to contribute substantially to the MI estimation error.

The fixed-margin expansion in Appendix~\ref{sec:fixed-margin-local-derivation}
relates MI, the pointwise-MI variance \(V(P)\), and the variance of the
sensitivities \(\tau^2(P)\) under the assumptions stated there. Close to
independence along that path, with fixed table dimensions and positive
cell probabilities, the expansion gives
\begin{equation}
 V(P)\approx 2I(P),
 \qquad \tau^2(P)\approx4V(P).
 \label{eq:near-independent-moments}
\end{equation}
Substituting these local approximations into
the component degrees-of-freedom formula gives
\begin{equation}
 \nu_{P,\mathrm{FO}}=\frac{2n_PV(P)^2}{\tau^2(P)}
 \approx\frac{n_PV(P)}{2}
 \approx n_PI(P).
 \label{eq:near-independent-df}
\end{equation}
At a fixed sample size, this population first-order value tends to zero
as positive population MI approaches zero. When \(\Vhat(P)\) and
\(\widehat\tau^2(P)\) estimate their population counterparts accurately
enough for the component formula to be reliable, \(\widehat\nu_P\) follows this
prediction: smaller MI gives fewer degrees of freedom and a larger cutoff
from the Student \(t\) distribution.

The empirical table, however, need not be independent even when the
population is independent. Random differences between its observed cell
proportions and the products of its observed margins produce nonzero
pointwise-MI values in occupied cells. Both \(\Vhat(P)\) and \(\widehat\tau^2(P)\) are
calculated from these values, so sampling-induced association enters the
estimated degrees of freedom. At a fixed sample size, these estimated
inputs need not track the shrinking population quantities accurately
enough for \(\widehat\nu_P\) to follow \(n_PI(P)\). Consequently, the
estimated degrees of freedom need not approach zero with population MI.
Appendix~\ref{sec:independence-boundary} retains the quadratic sampling
terms and derives their contribution to the estimated degrees of freedom.

\begin{samepage}
\subsection{Theoretical status and calculation requirements}

Equation~\eqref{eq:g-final} is an exact derivative for positive population
tables. Applying it to finite table changes gives the first-order
approximation to the variance-estimation error in
Equation~\eqref{eq:vhat-linearization}. The corresponding MI sampling-error
approximation was derived in Section~\ref{sec:first-order-mi-background}. Moment matching then
approximates each variance contribution by a variable with a scaled
chi-squared distribution, as in Equation~\eqref{eq:scaled-chi-working}.
An exact Student \(t\) distribution would require a numerator with a
standard normal distribution and a squared denominator with the scaled
chi-squared distribution whose mean is one. The numerator and squared denominator must be
independent. Here, the corrected MI
difference and its estimated standard error come from the same tables, so
an observation can change both quantities. Matching the variance moments
does not establish the required
distributional or independence conditions, and the Student \(t\) distribution
for \(T\) therefore remains a finite-sample approximation.

For fixed population tables satisfying the stated assumptions, including
positive \(V\) and \(\tau^2\), the estimated degrees of freedom grow as
the sample sizes tend to infinity, bringing the Student \(t\) distribution
closer to the standard normal distribution used by the Wald test. This large-sample behaviour holds the
populations fixed, whereas the near-independence comparison above reduces
population MI at a fixed sample size.
\par\nopagebreak[4]

Both tests need a finite corrected MI difference and a positive, finite
estimated standard error, while the Expanded Welch test also needs positive, finite component
and combined degrees of freedom. If a requirement is not met, the calculation
is unsuccessful and cannot be used to draw a test conclusion.
The experiment applies the numerical calculation requirements in
Appendix~\ref{sec:validity-rules} and assesses whether false-positive rates
are close to 5\%.
\par
\end{samepage}
